\chapter{Introduction}
\label{ch:introduction}

\section{Motivation}
\label{sec:motivation}

WebAssembly is a binary instruction format for a stack-based virtual machine. Its most mainstream use case is in browsers. Many languages can be compiled to it, including Kotlin.

There is a proposal to implement stack-switching in WebAssembly. The specification for stack-switching is finalized and its implementation is on the way.

The current implementation of Kotlin Coroutines for Wasm (with state-machines and continuation passing style for suspend function) is based on the implementation of suspend functions for Kotlin/JVM, which doesn't have primitives for continuation objects and suspension of functions.

The ultimate goal is to verify whether compiling Kotlin Coroutines using stack-switching improves binary size, complexity and performance for available Wasm runtimes with stack-switching support.

There is also an indirect benefit of preparing for eventual release of stack-switching in the mainstream Wasm runtimes.

\section{Thesis Structure}
\label{sec:structure}

The thesis is organized as follows: Chapter~\ref{ch:background} provides background information on WebAssembly stack-switching and Kotlin coroutines primitives. Chapter~\ref{ch:implementation} describes the implementation of Kotlin coroutines using the WebAssembly stack-switching proposal. Chapter~\ref{ch:evaluation} evaluates the implementation in terms of performance and binary size. Finally, Chapter~\ref{ch:conclusion} concludes the thesis and discusses future work.
